\par\Needspace{8\baselineskip}\bigskip
\pdfbookmark[3]{Exposition, 2}{pallas-II-article2}
\begin{center}2.\end{center}
\noindent\hspace*{\parindent}Substituant, dans l’équation 16, \(\mathrm{d}x^2+\mathrm{d}y^2+\mathrm{d}z^2=\mathrm{d}r^2+\frac{pm\mathrm{d}t^2}{rr}\mathpunct{\raisebox{.45ex}{,}}\) on trouve \(\frac{\mathrm{d}r^2}{\mathrm{d}t^2}=-\frac{m}{a}+\frac{2m}{r}-\frac{pm}{rr}\mathpunct{\raisebox{.45ex}{;}}\) ainsi dans les points de l’orbite, où \(r\) est un maximum \SourceLeaf{449}ou un minimum, on doit avoir \(\frac{m}{a}-\frac{2m}{r}+\frac{pm}{rr}=0\). Cela donne pour \(r\) deux valeurs, savoir \(r=a+\BareRoot{(aa-ap)}\) et \(r=a-\BareRoot{(aa-ap)}\), d’où il est évident, que \(2a\) est le grand axe, et \(\frac{\BareRoot{(aa-ap)}}{a}\mathpunct{\raisebox{.45ex}{,}}\) laquelle quantité nous ferons \(=e\), l’excentricité. On a donc
\[
\begin{gathered}
p=a(1-ee)\\[4pt]
\mathrm{d}t^2=\frac{arr\mathrm{d}r^2}{m(aaee-(a-r)^2)}\mathpunct{\cdot}
\end{gathered}
\]
\par Pour intégrer cette équation, nous ferons
\[
\begin{array}{rl}17)\quad&a-r=ae\cos E,\end{array}
\]
d’où il s’ensuit
\[
\mathrm{d}t^2=\frac{arr\mathrm{d}E^2}{m}\mathpunct{\raisebox{.45ex}{;}}
\]
donc supposant, ce qui est permis, que \(E\) aille toujours en croissant, il sera
\[
\mathrm{d}t=\BareRoot{\frac{a}{m}}\cdot r\mathrm{d}E=\BareRoot{\frac{a^3}{m}}\cdot(1-e\cos E)\mathrm{d}E.
\]
\par Ainsi en faisant
\[
\begin{array}{rl}
18)\quad&\BareRoot{\dfrac{m}{a^3}}=n\\[6pt]
19)\quad&E-e\sin E-nt=\epsilon,
\end{array}
\]
l’équation \(n\mathrm{d}t=(1-e\cos E)\mathrm{d}E\) montre que \(\epsilon\) sera une constante. On voit que \(E\) représente ce que les astronomes nomment l’anomalie excentrique, \(nt+\epsilon\) l’anomalie moyenne, \(n\) le mouvement moyen dans l’unité du tems, \(\epsilon\) l’époque de l’anomalie moyenne pour le tems zéro. De plus il est clair que dans le tems d’une révolution \(E\) doit croître de \(360^{0}\) ou de \(2\pi\), \(2\pi\) étant la périphérie du cercle dont le rayon est \(1\): ainsi nommant \(T\) ce tems, on aura
\[
nT=2\pi\qquad\text{ou}\qquad T=2\pi\BareRoot{\frac{a^3}{m}}\mathpunct{\raisebox{.45ex}{,}}
\]
ce qui revient au théorème de \textsc{Kepler}.
\par Par là on a encore la mesure des masses, lorsqu’on a choisi l’unité du tems et de l’espace, savoir
\[
m=\frac{4\pi\pi a^3}{TT}\mathpunct{\cdot}
\]
\par Ainsi choisissant pour l’une le jour solaire moyen et pour l’autre la \SourceLeaf{450}valeur de \(a\) dans l’orbite de la terre, supposant de plus \(T=365{,}2[563835]\), on trouve
\begin{center}
\textit{Somme des masses du soleil et de la terre}\par\nobreak
\mbox{. . . . \(0{,}000295913[0424]\)}
\end{center}
ou en supposant la première \(354710\) fois plus grande que la dernière
\begin{center}
\textit{Masse du soleil} . . . . \(0{,}000295912[2082]\),
\end{center}
dont le logarithme ordinaire est \(6{,}471\;1628\;828-10\).
